\chapter{2016年天津卷（文）}

\section{单选题}

\begin{problem}
已知集合 \( A=\{1,2,3\} \)，\( B=\{y\mid y=2x-1,\ x\in A\} \)，则 \( A\cap B=\)（\quad）

\choices
    {\(\{1,3\}\)}
    {\(\{1,2\}\)}
    {\(\{2,3\}\)}
    {\(\{1,2,3\}\)}
\end{problem}

\begin{answer}
A
\end{answer}

\begin{solution}
由 \(x\in A\) 可得 \(x=1\)，\(2\)，\(3\)，所以
\(B=\{2\times1-1,2\times2-1,2\times3-1\}=\{1,3,5\}\). 因此
\(A\cap B=\{1,3\}\)，选 A.
\end{solution}

\begin{problem}
甲，乙两人下棋，两人下成和棋的概率是 \( \frac12 \)，甲获胜的概率是 \( \frac13 \)，则甲不输的概率为（\quad）

\choices
    {\(\frac56\)}
    {\(\frac25\)}
    {\(\frac16\)}
    {\(\frac13\)}
\end{problem}

\begin{answer}
A
\end{answer}

\begin{solution}
甲不输包含甲获胜和两人和棋这两种互斥情况，所以甲不输的概率为
\(\dfrac13+\dfrac12=\dfrac56\)，选 A.
\end{solution}

\begin{problem}
将一个长方体沿相邻三个面的对角线截去一个棱锥，得到的几何体的正视图与俯视图如图所示，则该几何体的侧（左）视图为（\quad）

\bitmapfigure[width=0.42\linewidth]{img/2016/tianjin_liberal/q3_fig2.png}

\choices
    {\choicebitmap[width=0.15\paperwidth]{img/2016/tianjin_liberal/q3_optA.png}}
    {\choicebitmap[width=0.15\paperwidth]{img/2016/tianjin_liberal/q3_optB.png}}
    {\choicebitmap[width=0.15\paperwidth]{img/2016/tianjin_liberal/q3_optC.png}}
    {\choicebitmap[width=0.15\paperwidth]{img/2016/tianjin_liberal/q3_optD.png}}
\end{problem}

\begin{answer}
B
\end{answer}

\begin{solution}
正视图中的截面线段连接矩形的左上顶点和右下顶点，俯视图中的截面线段连接矩形的左下顶点和右上顶点. 将同一截面线段投影到左视图后，其两个端点分别落在左视图矩形的左上、右下位置；从左侧观察时该线段位于实体的后方，应画成虚线. 因此侧（左）视图为选项 B.
\end{solution}

\begin{problem}
已知双曲线 \( \dfrac{x^2}{a^2}-\dfrac{y^2}{b^2}=1 \)（\(a>0\)，\(b>0\)）的焦距为 \( 2\sqrt5 \)，且双曲线的一条渐近线与直线 \( 2x+y=0 \) 垂直，则双曲线的方程为（\quad）

\choices
    {\(\dfrac{x^2}{4}-y^2=1\)}
    {\(x^2-\dfrac{y^2}{4}=1\)}
    {\(\dfrac{3x^2}{20}-\dfrac{3y^2}{5}=1\)}
    {\(\dfrac{3x^2}{5}-\dfrac{3y^2}{20}=1\)}
\end{problem}

\begin{answer}
A
\end{answer}

\begin{solution}
双曲线的焦距为 \(2\sqrt5\)，故焦半距 \(c=\sqrt5\)，从而
\(a^2+b^2=c^2=5\). 双曲线的渐近线斜率为 \(\pm\dfrac ba\)，而直线
\(2x+y=0\) 的斜率为 \(-2\). 与该直线垂直的渐近线斜率为 \(\dfrac12\)，所以
\(\dfrac ba=\dfrac12\). 于是 \(b^2=\dfrac{a^2}{4}\)，代入
\(a^2+b^2=5\) 得 \(a^2=4\)，\(b^2=1\). 因此双曲线方程为
\(\dfrac{x^2}{4}-y^2=1\)，选 A.
\end{solution}

\begin{problem}
设 \(x>0\)，\(y\in \R\)，则“\( x>y \)”是“\( x>|y| \)”的（\quad）

\choices
    {充要条件}
    {充分不必要条件}
    {必要而不充分条件}
    {既不充分也不必要条件}
\end{problem}

\begin{answer}
C
\end{answer}

\begin{solution}
若 \(x>|y|\)，则 \(|y|\geqslant y\)，从而 \(x>y\)，所以“\(x>|y|\)”是“\(x>y\)”的充分条件.
反之不成立，例如 \(x=1\)，\(y=-2\) 时 \(x>y\)，但 \(x>|y|\) 不成立. 因此“\(x>y\)”是“\(x>|y|\)”的必要而不充分条件，选 C.
\end{solution}

\begin{problem}
已知 \( f(x) \) 是定义在 \( \R \) 上的偶函数，且在区间 \( \left( -\infty,0 \right) \) 上单调递增，若实数 \( a \) 满足 \( f\left( 2^{|a-1|} \right)>f\left( -\sqrt2 \right) \)，则 \( a \) 的取值范围是（\quad）

\choices
    {\(\left( -\infty,\frac12 \right)\)}
    {\(\left( -\infty,\frac12 \right)\cup\left( \frac32,+\infty \right)\)}
    {\(\left( \frac12,\frac32 \right)\)}
    {\(\left( \frac32,+\infty \right)\)}
\end{problem}

\begin{answer}
C
\end{answer}

\begin{solution}
令 \(u=2^{|a-1|}>0\). 由 \(f\) 为偶函数，\(f(u)=f(-u)\). 在
\((-\infty,0)\) 上 \(f\) 单调递增，因此
\[
f(u)>f(-\sqrt2)\iff f(-u)>f(-\sqrt2)\iff -u>-\sqrt2
\]
所以 \(u<\sqrt2\)，即
\(2^{|a-1|}<2^{1/2}\). 由于底数 \(2>1\)，有
\(|a-1|<\dfrac12\)，故
\(\dfrac12<a<\dfrac32\)，选 C.
\end{solution}

\begin{problem}
已知 \( \triangle ABC \) 是边长为 \( 1 \) 的等边三角形，点 \(D\)，\(E\) 分别是边 \(AB\)，\(BC\) 的中点，连接 \(DE\) 并延长到点 \(F\)，使得 \( DE=2EF \)，则 \( \overrightarrow{AF}\cdot\overrightarrow{BC} \) 的值为（\quad）

\choices
    {\(-\frac58\)}
    {\(\frac18\)}
    {\(\frac14\)}
    {\(\frac{11}{8}\)}
\end{problem}

\begin{answer}
B
\end{answer}

\begin{solution}
建立平面直角坐标系，取 \(A=(0,0)\)，\(B=(1,0)\)，
\(C=\left(\dfrac12,\dfrac{\sqrt3}{2}\right)\). 则
\(D=\left(\dfrac12,0\right)\)，\(E=\left(\dfrac34,\dfrac{\sqrt3}{4}\right)\)
因此
\(\overrightarrow{DE}=\left(\dfrac14,\dfrac{\sqrt3}{4}\right)\). 由 \(DE=2EF\) 且 \(DE\) 延长到 \(F\)，得
\(\overrightarrow{EF}=\dfrac12\overrightarrow{DE}\)，从而
\(F=\left(\dfrac78,\dfrac{3\sqrt3}{8}\right)\). 于是
\(\overrightarrow{AF}=\left(\dfrac78,\dfrac{3\sqrt3}{8}\right)\)，\(qquad \overrightarrow{BC}=\left(-\dfrac12,\dfrac{\sqrt3}{2}\right)\)
\[
\overrightarrow{AF}\cdot\overrightarrow{BC}
=-\dfrac{7}{16}+\dfrac{9}{16}=\dfrac18
\]
故选 B.
\end{solution}

\begin{problem}
已知函数 \( f(x)=\sin^2\frac{\omega x}{2}+\frac12\sin\omega x-\frac12 \)（\( \omega>0 \)，\( x\in \R \)），若 \( f(x) \) 在区间 \( \left( \pi,2\pi \right) \) 内没有零点，则 \( \omega \) 的取值范围是（\quad）

\choices
    {\(\left( 0,\frac18 \right]\)}
    {\(\left( 0,\frac14 \right]\cup\left[ \frac58,1 \right)\)}
    {\(\left( 0,\frac58 \right]\)}
    {\(\left( 0,\frac18 \right]\cup\left[ \frac14,\frac58 \right]\)}
\end{problem}

\begin{answer}
D
\end{answer}

\begin{solution}
利用 \(\sin^2\dfrac{\omega x}{2}=\dfrac{1-\cos\omega x}{2}\)，得
\[
f(x)=\dfrac{\sin\omega x-\cos\omega x}{2}
\]
所以 \(f(x)=0\) 等价于
\(\tan\omega x=1\)，即存在整数 \(k\) 使
\(\omega x=\dfrac\pi4+k\pi\). 令 \(t=\dfrac{\omega x}{\pi}\)，当
\(x\in(\pi,2\pi)\) 时 \(t\in(\omega,2\omega)\). 因此题意等价于开区间
\((\omega,2\omega)\) 不含形如 \(\dfrac14+k\) 的数.

当 \(0<\omega\leqslant\dfrac18\) 时，
\((\omega,2\omega)\subset(0,\dfrac14]\)，不含 \(\dfrac14\). 当
\(\dfrac18<\omega<\dfrac14\) 时，\(\dfrac14\in(\omega,2\omega)\).
当 \(\dfrac14\leqslant\omega\leqslant\dfrac58\) 时，
\((\omega,2\omega)\subset[\dfrac14,\dfrac54]\)，两端的可能零点只在开区间外，且区间内不含其他零点.
当 \(\dfrac58<\omega<\dfrac54\) 时，\(\dfrac54\in(\omega,2\omega)\). 当
\(\omega\geqslant\dfrac54\) 时，令 \(k=\lfloor\omega-\dfrac14\rfloor+1\)，则
\(\omega<\dfrac14+k\leqslant\omega+1<2\omega\)，所以
\(\dfrac14+k\in(\omega,2\omega)\).
因此
\(\omega\in\left(0,\dfrac18\right]\cup\left[\dfrac14,\dfrac58\right]\)，选 D.
\end{solution}

\section{填空题}

\begin{problem}
\(\i\) 是虚数单位，复数 \( z \) 满足 \( \left( 1+\i \right)z=2 \)，则 \( z \) 的实部为\fillinblank{}.
\end{problem}

\begin{answer}
1
\end{answer}

\begin{solution}
由
\[
z=\dfrac{2}{1+\i}=\dfrac{2(1-\i)}{(1+\i)(1-\i)}=1-\i
\]
可知 \(z\) 的实部为 \(1\).
\end{solution}

\begin{problem}
已知函数 \( f(x)=\left( 2x+1 \right)\e^x \)，\( f'(x) \) 为 \( f(x) \) 的导函数，则 \( f'(0) \) 的值为\fillinblank{}.
\end{problem}

\begin{answer}
3
\end{answer}

\begin{solution}
由乘积求导法则，
\(f'(x)=2\e^x+(2x+1)\e^x=(2x+3)\e^x\). 所以
\(f'(0)=3\e^0=3\).
\end{solution}

\begin{problem}
阅读如图所示的程序框图，运行相应的程序，则输出 \(S\) 的值为\fillinblank{}.

\bitmapfigure[width=0.38\linewidth]{img/2016/tianjin_liberal/q11_fig1.png}
\end{problem}

\begin{answer}
4
\end{answer}

\begin{solution}
按程序逐次运行：初始 \(S=4\)，\(n=1\). 第一次判断 \(S<6\)，所以
\(S=2S=8\)，再令 \(n=2\)；第二次判断 \(S\geqslant6\)，所以
\(S=S-6=2\)，再令 \(n=3\)；第三次判断 \(S<6\)，所以
\(S=2S=4\)，再令 \(n=4\). 此时 \(n>3\)，程序输出 \(S=4\).
\end{solution}

\begin{problem}
已知圆 \(C\) 的圆心在 \( x \) 轴正半轴上，点 \(M\left( 0,\sqrt5 \right) \) 在圆 \(C\) 上，且圆心到直线 \( 2x-y=0 \) 的距离为 \( \frac{4\sqrt5}{5} \)，则圆 \(C\) 的方程为\fillinblank{}.
\end{problem}

\begin{answer}
\( (x-2)^2+y^2=9 \)
\end{answer}

\begin{solution}
设圆心为 \(C=(r,0)\)，其中 \(r>0\). 由 \(M=(0,\sqrt5)\) 在圆上，圆的半径平方为
\(r^2+5\). 又圆心到直线 \(2x-y=0\) 的距离为
\[
\dfrac{|2r|}{\sqrt{2^2+(-1)^2}}=\dfrac{2r}{\sqrt5}=\dfrac{4\sqrt5}{5}=\dfrac4{\sqrt5}
\]
所以 \(r=2\)，半径平方为 \(2^2+5=9\). 因此圆 \(C\) 的方程为
\((x-2)^2+y^2=9\).
\end{solution}

\begin{problem}
如图，\(AB\) 是圆的直径，弦 \(CD\) 与 \(AB\) 相交于点 \(E\)，\( BE=2AE=2 \)，\( BD=ED \)，则线段 \(CE\) 的长为\fillinblank{}.

\bitmapfigure[width=0.33\linewidth]{img/2016/tianjin_liberal/q13_fig1.png}
\end{problem}

\begin{answer}
\(\frac{2\sqrt3}{3}\)
\end{answer}

\begin{solution}
由 \(BE=2AE=2\)，得 \(AE=1\)，\(BE=2\). 取 \(E\) 为原点，令
\(A=(-1,0)\)，\(B=(2,0)\). 因 \(AB\) 是直径，圆的方程为
\[
(x+1)(x-2)+y^2=0
\]
即 \(x^2-x+y^2-2=0\). 又 \(BD=ED\)，所以 \(D\) 在 \(BE\) 的垂直平分线上，故 \(D=(1,y)\). 代入圆方程得 \(y^2=2\)，从而
\(ED=\sqrt{1^2+2}=\sqrt3\).

由相交弦定理，
\(EA\cdot EB=EC\cdot ED\)，所以
\[
EC=\dfrac{1\times2}{\sqrt3}=\dfrac{2\sqrt3}{3}
\]
故线段 \(CE\) 的长为 \(\dfrac{2\sqrt3}{3}\).
\end{solution}

\begin{problem}
已知函数
\( f(x)=\begin{cases}
x^2+\left( 4a-3 \right)x+3a, & x<0,\\
\log_a\left( x+1 \right)+1, & x\ge 0
\end{cases} \)
（\( a>0 \)，且 \( a\ne 1 \)）在 \( \R \) 上单调递减，且关于 \( x \) 的方程 \( |f(x)|=2-\frac{x}{3} \) 恰有两个不相等的实数解，则 \( a \) 的取值范围是\fillinblank{}.
\end{problem}

\begin{answer}
\(\left[\frac13,\frac23\right)\)
\end{answer}

\begin{solution}
记 \(p(x)=x^2+(4a-3)x+3a\)，\(q(x)=\log_a(x+1)+1\). 在 \(x<0\) 时，
\(p'(x)=2x+4a-3\). 要使 \(p\) 在 \((-\infty,0)\) 上递减，需
\(a\leqslant\dfrac34\)；在 \(x\geqslant0\) 时，\(q\) 递减要求 \(0<a<1\). 此外，为保证跨过 \(0\) 仍递减，需
\(\lim_{x\to0^-}p(x)=3a\geqslant q(0)=1\)，即 \(a\geqslant\dfrac13\). 反之在
\(\dfrac13\leqslant a\leqslant\dfrac34\) 时，上述条件均成立，所以
\[
\dfrac13\leqslant a\leqslant\dfrac34
\]

在这个范围内，若 \(x<0\)，则 \(p(x)>0\)，方程化为
\[
h(x)=x^2+\left(4a-\dfrac83\right)x+3a-2=0
\]
当 \(a<\dfrac23\) 时，\(h(0)<0\)，且二次项系数为正、常数项为负，所以恰有一个负根；当 \(a=\dfrac23\) 时 \(h(x)=x^2\)，没有负根；当 \(a>\dfrac23\) 时，
\[
\Delta=\dfrac49(36a^2-75a+34)
=\dfrac49\cdot36\left(a-\dfrac23\right)\left(a-\dfrac{17}{12}\right)<0
\]
没有实根. 因此负半轴贡献一个解当且仅当
\(\dfrac13\leqslant a<\dfrac23\).

再看 \(x\geqslant0\). 当 \(q(x)\geqslant0\) 时，由 \(0<a<1\) 得
\(0\leqslant x\leqslant a^{-1}-1\leqslant2\)，从而
\(|q(x)|=q(x)\leqslant1<2-\dfrac{x}{3}\)，没有解. 当 \(q(x)<0\) 时，令
\[
G(x)=q(x)+2-\dfrac{x}{3}=\log_a(x+1)+3-\dfrac{x}{3}
\]
因为 \(a<1\)，有
\(G'(x)=\dfrac{1}{(x+1)\ln a}-\dfrac13<0\). 在
\(x_0=a^{-1}-1\) 处 \(G(x_0)=2-\dfrac{x_0}{3}>0\). 又因
\(a\geqslant\dfrac13\)，有 \(a^{-1}\leqslant3<7\)，且 \(0<a<1\)，所以
\(\log_a7<\log_a(a^{-1})=-1\)，即 \(G(6)=\log_a7+1<0\). 因此由连续性和严格递减性，正半轴恰有一个解.

综上，原方程恰有两个不相等的实数解当且仅当
\(\boxed{a\in\left[\dfrac13,\dfrac23\right)}\).
\end{solution}

\section{解答题}

\begin{problem}
在 \( \triangle ABC \) 中，内角 \(A\)，\(B\)，\(C\) 所对的边分别为 \( a \)，\( b \)，\( c \)，已知 \( a\sin 2B=\sqrt3 b\sin A \).

\begin{enumerate}
\item 求 \(B\)；
\item 已知 \( \cos A=\dfrac13 \)，求 \( \sin C \) 的值.
\end{enumerate}
\end{problem}

\begin{solution}
由正弦定理，\(\dfrac{a}{\sin A}=\dfrac{b}{\sin B}\)，所以
\(a=\dfrac{b\sin A}{\sin B}\). 代入已知等式并利用
\(\sin 2B=2\sin B\cos B\)，得
\[
\dfrac{b\sin A}{\sin B}\cdot2\sin B\cos B
=\sqrt3\,b\sin A
\]
由于 \(b>0\) 且 \(0<A<\pi\)，可约去 \(b\sin A\)，得到
\(2\cos B=\sqrt3\). 三角形内角满足 \(0<B<\pi\)，故
\(B=\dfrac{\pi}{6}\).

又已知 \(\cos A=\dfrac13\)，且 \(0<A<\pi\)，所以
\(\sin A=\sqrt{1-\cos^2A}=\dfrac{2\sqrt2}{3}\). 由
\(C=\pi-A-B\)，有
\[
\sin C=\sin(A+B)=\sin A\cos B+\cos A\sin B
=\dfrac{2\sqrt2}{3}\cdot\dfrac{\sqrt3}{2}
+\dfrac13\cdot\dfrac12
=\dfrac{2\sqrt6+1}{6}
\]
因此 \(B=\dfrac{\pi}{6}\)，\(\sin C=\dfrac{2\sqrt6+1}{6}\).
\end{solution}

\begin{problem}
某化工厂生产甲，乙两种混合肥料，需要 \(A\)，\(B\)，\(C\) 三种主要原料，生产 1 车皮甲种肥料和生产 1 车皮乙种肥料所需三种原料的吨数如表所示：

\begin{center}
\begin{tabular}{|c|c|c|c|}
\hline
\begin{tabular}{c}原料\\肥料\end{tabular} & \(A\) & \(B\) & \(C\) \\
\hline
甲 & 4 & 8 & 3 \\
乙 & 5 & 5 & 10 \\
\hline
\end{tabular}
\end{center}

现有 \(A\) 种原料 200 吨，\(B\) 种原料 360 吨，\(C\) 种原料 300 吨，在此基础上生产甲，乙两种肥料. 已知生产 1 车皮甲种肥料，产生的利润为 2 万元；生产 1 车皮乙种肥料，产生的利润为 3 万元. 分别用 \( x \)，\( y \) 表示计划生产甲，乙两种肥料的车皮数.

\begin{enumerate}
\item 用 \( x \)，\( y \) 列出满足生产条件的数学关系式，并画出相应的平面区域；
\item 问分别生产甲，乙两种肥料各多少车皮，能够产生最大的利润？并求出此最大利润.
\end{enumerate}
\end{problem}

\begin{solution}
\begin{enumerate}
\item 由三种原料的用量限制以及 \(x\)，\(y\geqslant0\)，可得可行域满足
\[
\begin{cases}
4x+5y\leqslant200,\\
8x+5y\leqslant360,\\
3x+10y\leqslant300,\\
x\geqslant0,\\
y\geqslant0.
\end{cases}
\]
三条资源约束直线在坐标轴上的截距分别为
\((50,0)\)，\((0,40)\)、\((45,0)\)，\((0,72)\)、\((100,0)\)，\((0,30)\)，所以坐标轴上的可行顶点为
\((45,0)\) 和 \((0,30)\). 三条直线两两相交时，第一、二条给出
\(4x+5y=200\)，\(8x+5y=360\)，相减得 \(x=40\)，代回得 \(y=8\)，即
\((40,8)\)；第一、三条给出
\(4x+5y=200\)，\(3x+10y=300\)，将第一式乘 \(2\) 后相减得 \(x=20\)，代回得 \(y=24\)，即
\((20,24)\). 第二、三条的交点为
\(\left(\dfrac{420}{13},\dfrac{264}{13}\right)\)，但在该点
\(4x+5y=\dfrac{3000}{13}>200\)，不满足第一条约束. 因此可行域的顶点依次为
\((0,0)\)，\((45,0)\)，\((40,8)\)，\((20,24)\)，\((0,30)\)
因此相应平面区域是第一象限内由这五个顶点围成的五边形，边界依次落在
\(y=0\)、\(8x+5y=360\)、\(4x+5y=200\)、\(3x+10y=300\)、\(x=0\) 上.

\item 设利润为 \(P\) 万元，则
\(P=2x+3y\). 在可行域的各顶点处计算：
\(P(0,0)=0\)，\(P(45,0)=90\)，\(P(40,8)=104\)，\(P(20,24)=112\)，\(P(0,30)=90\)
线性目标函数在多边形可行域上的最大值必在顶点取得，所以最大利润为
\(112\) 万元，此时应生产甲 \(20\) 车皮、乙 \(24\) 车皮.
\end{enumerate}
\end{solution}

\begin{problem}
如图，四边形 \(ABCD\) 是平行四边形，平面 \( AED \perp \) 平面 \(ABCD\)，\( EF\parallel AB \)，\( AB=2 \)，\( DE=3 \)，\( BC=EF=1 \)，\( AE=\sqrt6 \)，\( \angle BAD=60^\circ \)，\(G\) 为 \(BC\) 的中点.

\begin{enumerate}
\item 求证：\( FG\parallel \) 平面 \(BED\)；
\item 求证：平面 \( BED \perp \) 平面 \(AED\)；
\item 求直线 \(EF\) 与平面 \(BED\) 所成角的正弦值.
\end{enumerate}

\FigureLayoutDeclare{img/2016/tianjin_liberal/q17_fig1.png}{12.0cm}{45bp 78bp 81bp 32bp}
\bitmapfigure[float=inline,width=0.55\linewidth]{img/2016/tianjin_liberal/q17_fig1.png}
\FloatBarrier
\end{problem}

\begin{solution}
\begin{enumerate}
\item 建立直角坐标系，使底面 \(ABCD\) 为 \(z=0\) 平面，取
\(A=(0,0,0)\)，\(B=(2,0,0)\)，\(D=\left(\dfrac12,\dfrac{\sqrt3}{2},0\right)\)，\(C=\left(\dfrac52,\dfrac{\sqrt3}{2},0\right)\)
设 \(E\) 在平面 \(AED\) 内的坐标为
\(\left(\dfrac t2,\dfrac{\sqrt3t}{2},z\right)\). 由
\(AE=\sqrt6\)，\(DE=3\)，\(AD=1\)，有
\(t^2+z^2=6\)，\((t-1)^2+z^2=9\)
两式相减得 \(t=-1\)，再由第一式得 \(z^2=5\). 按图取
\(E=\left(-\dfrac12,-\dfrac{\sqrt3}{2},\sqrt5\right)\). 由于
\(EF\parallel AB\)、\(EF=1\)，取
\(F=\left(\dfrac12,-\dfrac{\sqrt3}{2},\sqrt5\right)\)，且
\[
G=\dfrac{B+C}{2}
=\left(\dfrac94,\dfrac{\sqrt3}{4},0\right)
\]
于是
\(\overrightarrow{BD}=\left(-\dfrac32,\dfrac{\sqrt3}{2},0\right)\)，\(\overrightarrow{BE}=\left(-\dfrac52,-\dfrac{\sqrt3}{2},\sqrt5\right)\)
\[
\overrightarrow{FG}
=\left(\dfrac74,\dfrac{3\sqrt3}{4},-\sqrt5\right)
=\dfrac12\overrightarrow{BD}-\overrightarrow{BE}
\]
所以 \(\overrightarrow{FG}\) 与平面 \(BED\) 内两条相交直线的方向向量共面，即
\(FG\) 的方向平行于平面 \(BED\). 又
\[
\left(\overrightarrow{BD}\times\overrightarrow{BE}\right)
\cdot\overrightarrow{BF}=\dfrac{\sqrt{15}}{2}\ne0
\]
故 \(F\notin\) 平面 \(BED\)，从而 \(FG\parallel\) 平面 \(BED\).

\item 由上述坐标，
\(\overrightarrow{AD}=\left(\dfrac12,\dfrac{\sqrt3}{2},0\right)\)，\(\overrightarrow{AE}=\left(-\dfrac12,-\dfrac{\sqrt3}{2},\sqrt5\right)\)
直接计算得
\(\overrightarrow{BD}\cdot\overrightarrow{AD}=0\)，\(\overrightarrow{BD}\cdot\overrightarrow{AE}=0\)
直线 \(AD\)，\(AE\) 是平面 \(AED\) 内相交的两条直线，所以
\(BD\perp\) 平面 \(AED\). 又 \(BD\subset\) 平面 \(BED\)，因此
\(\text{平面 }BED\perp\text{平面 }AED\).

\item 设直线 \(EF\) 与平面 \(BED\) 所成的角为 \(\theta\). 取
\(H\) 为 \(A\) 到 \(DE\) 的垂足. 由于两平面交于 \(DE\)，且
\(AED\perp BED\)，有 \(AH\perp\) 平面 \(BED\). 又 \(EF\parallel AB\)，所以
\(\theta\) 等于直线 \(AB\) 与平面 \(BED\) 所成的角，即
\(\theta=\angle ABH\).

在三角形 \(ADE\) 中，由余弦定理，
\[
\cos\angle ADE
=\dfrac{AD^2+DE^2-AE^2}{2AD\cdot DE}
=\dfrac{1+9-6}{2\times1\times3}
=\dfrac23
\]
所以 \(\sin\angle ADE=\dfrac{\sqrt5}{3}\)，从而
\(AH=AD\sin\angle ADE=\dfrac{\sqrt5}{3}\). 在直角三角形
\(AHB\) 中，\(AB=2\)，故
\[
\sin\theta=\sin\angle ABH=\dfrac{AH}{AB}
=\dfrac{\sqrt5}{6}
\]
\end{enumerate}
\end{solution}

\begin{problem}
已知 \( \{a_n\} \) 是等比数列，前 \( n \) 项和为 \( S_n \)（\( n\in \N^* \)），且 \( \dfrac{1}{a_1}-\dfrac{1}{a_2}=\dfrac{2}{a_3} \)，\( S_6=63 \).

\begin{enumerate}
\item 求 \( \{a_n\} \) 的通项公式；
\item 若对任意的 \( n\in \N^* \)，\( b_n \) 是 \( \log_2 a_n \) 和 \( \log_2 a_{n+1} \) 的等差中项，求数列 \( \{(-1)^n b_n^2\} \) 的前 \( 2n \) 项和.
\end{enumerate}
\end{problem}

\begin{solution}
\begin{enumerate}
\item 设等比数列的首项为 \(a_1\)，公比为 \(q\). 则
\(a_2=a_1q\)，\(a_3=a_1q^2\)，由已知
\[
\dfrac1{a_1}-\dfrac1{a_1q}=\dfrac2{a_1q^2}
\]
乘以 \(a_1q^2\) 得 \(q^2-q=2\)，即
\((q-2)(q+1)=0\)，所以 \(q=2\) 或 \(q=-1\). 若 \(q=-1\)，则
\(S_6=0\)，与 \(S_6=63\) 矛盾，故 \(q=2\). 于是
\[
S_6=a_1(1+2+\cdots+2^5)=63a_1=63
\]
所以 \(a_1=1\)，从而
\(a_n=2^{n-1}\).

\item 因为 \(b_n\) 是 \(\log_2a_n\) 与 \(\log_2a_{n+1}\) 的等差中项，
\[
b_n=\dfrac{\log_2a_n+\log_2a_{n+1}}2
=\dfrac{n-1+n}{2}=n-\dfrac12
\]
设所求前 \(2n\) 项和为 \(T_{2n}\)，两项一组配对，得
\[
T_{2n}=\sum_{j=1}^{n}\left[-\left(2j-\dfrac32\right)^2
+\left(2j-\dfrac12\right)^2\right]
=\sum_{j=1}^{n}(4j-2)=2n^2
\]
故前 \(2n\) 项和为 \(2n^2\).
\end{enumerate}
\end{solution}

\begin{problem}
设椭圆 \( \dfrac{x^2}{a^2}+\dfrac{y^2}{3}=1 \)（\( a>\sqrt3 \)）的右焦点为 \(F\)，右顶点为 \(A\)，已知 \( \dfrac1{|OF|}+\dfrac1{|OA|}=\dfrac{3e}{|FA|} \)，其中 \(O\) 为原点，\( e \) 为椭圆的离心率.

\begin{enumerate}
\item 求椭圆的方程；
\item 设过点 \(A\) 的直线 \( l \) 与椭圆交于 \(B\)（\(B\) 不在 \( x \) 轴上），垂直于 \( l \) 的直线与 \( l \) 交于点 \(M\)，与 \( y \) 轴交于点 \(H\)，若 \( BF\perp HF \)，且 \( \angle MOA=\angle MAO \)，求直线 \( l \) 的斜率.
\end{enumerate}
\end{problem}

\begin{solution}
\begin{enumerate}
\item 设椭圆的半长轴为 \(a\)，半短轴为 \(\sqrt3\)，则
\(c=\sqrt{a^2-3}\)，且 \(e=\dfrac ca\). 由
\(OF=c\)，\(OA=a\)，\(FA=a-c\)，已知条件化为
\[
\dfrac1c+\dfrac1a=\dfrac{3c}{a(a-c)}
\]
两边同乘 \(ac(a-c)\)，得
\(a^2-c^2=3c^2\)，即 \(a^2=4c^2\). 结合
\(a^2-c^2=3\)，得 \(c^2=1\)，\(a^2=4\). 所以椭圆方程为
\[
\dfrac{x^2}{4}+\dfrac{y^2}{3}=1
\]

\item 设直线 \(l\) 的斜率为 \(k\). 因 \(B\) 不在 \(x\) 轴上，\(k\ne0\)，且
\(l:y=k(x-2)\). 将其代入椭圆方程，除去对应于 \(A\) 的根 \(x=2\)，得到
\(x_B=\dfrac{8k^2-6}{3+4k^2}\)，\(y_B=-\dfrac{12k}{3+4k^2}\)
直线 \(l\) 的垂线过 \(M\)，且 \(\angle MOA=\angle MAO\). 在三角形
\(MOA\) 中，等角所对的边相等，即 \(OM=AM\). 若
\(M=(x_M,y_M)\)，则
\[
x_M^2+y_M^2=(x_M-2)^2+y_M^2
\]
所以 \(x_M=1\)，又 \(M\in l\)，故 \(M=(1,-k)\). 过 \(M\) 的垂线斜率为
\(-\dfrac1k\)，因此它与 \(y\) 轴的交点为
\[
H=\left(0,\dfrac1k-k\right)
\]
由 \(BF\perp HF\)，且 \(F=(1,0)\)，有
\[
(x_B-1,y_B)\cdot\left(-1,\dfrac1k-k\right)=0
\]
代入 \(x_B\)，\(y_B\) 得
\[
\dfrac{4k^2-9}{3+4k^2}
=\dfrac{12(k^2-1)}{3+4k^2}
\]
于是 \(12k^2-12=4k^2-9\)，即 \(k^2=\dfrac38\). 因此
\[
k=\pm\sqrt{\dfrac38}=\pm\dfrac{\sqrt6}{4}
\]
\end{enumerate}
\end{solution}

\begin{problem}
设函数 \( f(x)=x^3-ax-b \)，\( x\in \R \)，其中 \(a\)，\(b\in \R\).

\begin{enumerate}
\item 求 \( f(x) \) 的单调区间；
\item 若 \( f(x) \) 存在极值点 \( x_0 \)，且 \( f(x_1)=f(x_0) \)，其中 \( x_1\ne x_0 \)，求证：\( x_1+2x_0=0 \)；
\item 设 \( a>0 \)，函数 \( g(x)=|f(x)| \)，求证：\( g(x) \) 在区间 \( [-1,1] \) 上的最大值不小于 \( \dfrac14 \).
\end{enumerate}
\end{problem}

\begin{solution}
\begin{enumerate}
\item \(f'(x)=3x^2-a\). 当 \(a\leqslant0\) 时，\(f'(x)\geqslant0\)，所以
\(f(x)\) 在 \(\R\) 上单调递增. 当 \(a>0\) 时，令
\(r=\sqrt{\dfrac a3}=\dfrac{\sqrt{3a}}3\)，则
\(f'(x)>0\)，\(\left(x\in(-\infty,-r)\cup(r,+\infty)\right)\)，\(f'(x)<0\)，\(\left(x\in(-r,r)\right)\)
所以 \(f\) 在 \((-\infty,-r)\) 和 \((r,+\infty)\) 上单调递增，在
\((-r,r)\) 上单调递减.

\item \(x_0\) 是极值点，所以 \(f'(x_0)=0\)，即
\(a=3x_0^2\). 由 \(f(x_1)=f(x_0)\)，得
\[
(x_1-x_0)(x_1^2+x_1x_0+x_0^2-a)=0
\]
因 \(x_1\ne x_0\)，所以
\(x_1^2+x_1x_0+x_0^2=3x_0^2\)，即
\[
(x_1-x_0)(x_1+2x_0)=0
\]
再次利用 \(x_1\ne x_0\)，得到 \(x_1+2x_0=0\).

\item 对任意 \(x\in[0,1]\)，有
\[
f(x)-f(-x)=2(x^3-ax)
\]
设 \(M=\max_{-1\leqslant t\leqslant1}|f(t)|\)，则
\(2M\geqslant|f(x)-f(-x)|\)，从而
\[
M\geqslant|x^3-ax|
\]
分情况取 \(x\)：

当 \(0<a\leqslant\dfrac34\) 时，取 \(x=1\)，则
\[
|x^3-ax|=1-a\geqslant\dfrac14
\]
当 \(\dfrac34<a<\dfrac54\) 时，取
\(x=\sqrt{\dfrac a3}\in(0,1)\)，则
\[
|x^3-ax|
=\left|\dfrac a3\sqrt{\dfrac a3}-a\sqrt{\dfrac a3}\right|
=\dfrac{2a^{3/2}}{3\sqrt3}
>\dfrac14
\]
当 \(a\geqslant\dfrac54\) 时，仍取 \(x=1\)，得到
\[
|x^3-ax|=a-1\geqslant\dfrac14
\]
因此无论 \(a>0\) 取何值，都有
\[
\max_{-1\leqslant x\leqslant1}g(x)=M\geqslant\dfrac14
\]
\end{enumerate}
\end{solution}
